\documentclass[parskip=full,11pt]{scrartcl}
%\usepackage{pdfpages}
\usepackage[utf8]{inputenc}
\usepackage{amssymb}
\usepackage[T1]{fontenc}
\usepackage[german]{babel}
\usepackage[yyyymmdd]{datetime} 
\usepackage{hyperref}
\usepackage[toc, nonumberlist, automake]{} %added automake option
\usepackage{csquotes}
\usepackage{graphicx}
\usepackage{longtable}
\usepackage{multirow}
\usepackage{pbox}
\usepackage{array}
\newcolumntype{P}[1]{>{\centering\arraybackslash}p{#1}}
\hypersetup{
 		pdftitle={VForWater-Implementierung},
 }
\usepackage{fancyhdr}%<-------------to control headers and footers
\usepackage[a4paper,margin=1in,footskip=.25in]{geometry}
\fancyhf{}
\fancyfoot[C]{\thepage} %<----to get page number below text
\pagestyle{fancy} %<-------the page style itself
 
\title{Qualitätssicherung}
\subtitle{Autorisierungsmanagement für eine virtuelle Forschungsumgebung für Geodaten}
\author{Alex\\Anastasia\\Atanas\\Dannie\\ Houra\\Sonya\\}
\date{11.01.18}
 % define custom lists
\usepackage{enumitem}


\usepackage{linegoal,listings}
\newsavebox{\mylisting}
\makeatletter
\newcommand{\lstInline}[2][,]{%
	\begingroup%
	\lstset{#1}% Set any keys locally
	\begin{lrbox}{\mylisting}\lstinline!#2!\end{lrbox}% Store listing in \mylisting
	\setlength{\@tempdima}{\linegoal}% Space left on line.
	\ifdim\wd\mylisting>\@tempdima\hfill\\\fi% Insert line break
	\lstinline!#2!% Reset listing
	\endgroup%
}
\makeatother
\setlength{\parindent}{0pt}% Just for this example

\lstset{basicstyle=\footnotesize\ttfamily,breaklines=true}
\lstset{framextopmargin=50pt,frame=bottomline,showstringspaces=false,upquote=true}


 
\begin{document}
 
 \begin{titlepage}
 	
 	\begin{center}
 	\includegraphics[width=0.5\linewidth]{res/KITLogo.png}\\
 	\vspace{2cm}
 	{\scshape\LARGE\bfseries Qualitätssicherung \par}
 	\vspace{0.5cm}
 	{\scshape\Large Praxis der Softwareentwicklung\\}
 	\vspace{1cm}
 	{\scshape\Large Wintersemester 17/18\\}
 	\vspace{2cm}
 	{\huge\bfseries Autorisierungsmanagement für eine virtuelle Forschungsumgebung für Geodaten\par}
 	\vspace{2cm}
 	\vfill
 	{\bfseries {\Large Autoren}:\par}
 	{\Large Bachvarov, Aleksandar }\\
 	{\Large Dimitrov, Atanas }\\
 	{\Large Mortazavi Moshkenan, Houraalsadat }\\
 	{\Large Sakly, Khalil }\\
 	{\Large Slobodyanik, Anastasia }\\
 	{\Large Voneva, Sonya}\\
 	\vfill
 	{\large 14.03.18 \par}
 	\end{center}
 \end{titlepage}
 
 \tableofcontents
 \newpage
 \section{Einleitung}
Dieses Dokument beschreibt ...

 \newpage
 \section{Qualität}
 Wie wird Qualität gemessen.
 
 \newpage
\section{Komponenten-Test}
Views unit Tests

\begin{longtable}[c]{|p{0.4cm}|P{2.5cm}|p{3.5cm}|p{5cm}|p{3cm}|}
\caption{My caption}
\label{my-label}\\
\hline
\textbf{\#} & \textbf{Testcase}&\textbf{Methoden}& \textbf{Beschreibung} & \textbf{Ergebnis} \\ \hline
\endfirsthead
%
\endhead
%
\multirow{2}{*}{}1 & \multirow{2}{*}{} TestHomeView & test-not-logged-in & Getestet, ob die Homeseite gezeigt wird, wenn man nicht eingeloggt ist. Das wird durch den zurückgegebene HTTP-Statuscode gemacht.& HTTP Status Code 302: Der Benutzer wird zur Loginseite weitergeleitet  \\ \cline{3-5} &   & test-normal & Getestet, ob die Homeseite gezeigt wird, wenn man eingeloggt ist.  & HTPP Status Code 200: OK \\ \hline

\multirow{3}{*}{} 2& \multirow{3}{*}{} TestResourceManager&  test-not-logged-in & Getestet, ob die ResourceManagerseite gezeigt wird, wenn man nicht eingeloggt ist. & HTTP Status Code 302: Der Benutzer wird zur Loginseite weitergeleitet \\ \cline{3-5} & & test-logged-in-no-admin & Getestet, ob die ResourceManagerseite gezeigt wird, wenn man eingeloggt ist, aber zwar nicht als Admin. & HTTP Status Code 302: Der Benutzer wird zur Adminloginseite weitergeleitet \\ \cline{3-5} & & test-normal & Getestet, ob die ResourceManagerseite gezeigt wird, wenn man als Admin  eingeloggt ist. & HTPP Status Code 200: OK  \\ \hline

\multirow{3}{*}{} 3& \multirow{3}{*}{} TestUserManager&  test-not-logged-in & Getestet, ob die UserManagerseite gezeigt wird, wenn man nicht eingeloggt ist. & HTTP Status Code 302: Der Benutzer wird zur Loginseite weitergeleitet  \\ \cline{3-5} & & test-logged-in-no-admin & Getestet, ob die UserManagerseite gezeigt wird, wenn man eingeloggt ist, aber zwar nicht als Admin. & HTTP Status Code 302: Der Benutzer wird zur Adminloginseite weitergeleitet \\ \cline{3-5} & & test-normal & Getestet, ob die UserManagerseite gezeigt wird, wenn man als Admin  eingeloggt ist. & HTPP Status Code 200: OK \\ \hline

\multirow{5}{*}{} 4& \multirow{5}{*}{} TestProfileView& test-not-logged-in & Getestet, ob die Profileseite gezeigt wird, wenn man nicht eingeloggt ist.& HTTP Status Code 302: Der Benutzer wird zur Loginseite weitergeleitet  \\ \cline{3-5} &   & test-normal & Getestet, ob die Profileseite gezeigt wird, wenn man eingeloggt ist.  & HTPP Status Code 200: OK \\ \cline{3-5} 
                  &                   & test-pagination-user & Getestet, ob der Benutzer die zwei Zugriffsrequesten sieht, die ihm gesendet wurden. & Anzahl der gezeigten Requeste auf der Seite - 2 \\ \cline{3-5} 
                  &                   & test-pagination-admin-page-1 & Getestet, ob der Admin  vier von den fünf Requesten, die ihm gesendet wurden,  auf der ersten Seite sieht. & Anzahl der gezeigten Requeste auf der ersten Seite - 4 \\ \cline{3-5} 
                  &                   & test-pagination-admin-page-2 & Getestet, ob der Admin  das Fünfte von den fünf Requesten, die ihm gesendet wurden,  auf der zweiten Seite sieht.& Anzahl der gezeigten Requeste auf der zweiten Seite - 1 \\  \hline
   
                  
\multirow{3}{*}{} 5& \multirow{3}{*}{} TestMyResourcesView& test-not-logged-in & Getestet, ob die MyResourcesseite gezeigt wird, wenn man nicht eingeloggt ist.& HTTP Status Code 302: Der Benutzer wird zur Loginseite weitergeleitet \\ \cline{3-5} &   & test-normal & Getestet, ob die MyResourcesseite gezeigt wird, wenn man eingeloggt ist.  & HTPP Status Code 200: OK \\ \cline{3-5}& & test-resources-shown & Getestet, ob der Benutzer seine zwei Ressourcen auf der MyResourcesseite  sieht. &  Anzahl der gezeigten Ressourcen auf der Seite - 2  \\ \hline


\multirow{6}{*}{}6 & \multirow{6}{*}{} TestSendDeletionRequestView& test-not-logged-in & Getestet, ob man ein Löschrequest senden kann, wenn man nicht eingeloggt ist.& HTTP Status Code 302: Der Benutzer wird zur Loginseite weitergeleitet \\ \cline{3-5} 
                  &                   & test-not-existing-resource &  Getestet, man ein Löschrequest für eine Ressource senden kann, wenn die Ressource nicht existiert  &  HTTP Status Code 302: Der Benutzer wird zur MyResourcesseite weitergeleitet  \\ \cline{3-5} 
                  &                   & test-not-owner &Getestet, ob ein Benutzer, der die Ressource nicht besitzt, ein Löschrequest senden kann. & HTTP Status Code 403: Der Benutzer wird zur PermissionDeniedseite weitergeleitet    \\ \cline{3-5} 
                  &                   & test-staff-user &Getestet, ob der Admin ein Löschrequest senden kann.  &HTTP Status Code 302: Der Admin wird zur MyResourcesseite weitergeleitet \\ \cline{3-5} 
                  &                   & test-deletion-request-exists  & Getestet, was passiert, wenn ein Löschrequest für diese Ressource schon existiert.  & HTTP Status Code 302: Der Benutzer wird zur MyResourcesseite weitergeleitet   \\ \cline{3-5} 
                  &                  & test-normal  & Getestet, was es im normalen Fall passiert, wenn man eine Löschrequest sendet. & HTTP Status Code 302: Der Benutzer wird zur MyResourcesseite weitergeleitet   \\ \hline
                  
                  
\multirow{6}{*}{}7& \multirow{6}{*}{} TestCancelDeletionRequestView& test-not-logged-in & Getestet, ob man ein Löschrequest ablehnen kann, wenn man nicht eingeloggt ist.& HTTP Status Code 302: Der Benutzer wird zur Loginseite weitergeleitet \\ \cline{3-5} 
                  &                   &  &  &    \\ \cline{3-5} 
                  &                   & test-not-owner &Getestet, ob ein Benutzer, der die Ressource nicht besitzt, ein Löschrequest ablehnen kann. & 403   \\ \cline{3-5} 
                  &                   & test-staff-user &Getestet, ob der Admin ein Löschrequest ablehnen kann(Admin kann keine Löschrequeste senden).  &HTTP Status Code 302: Der Benutzer wird zur Loginseite weitergeleitet  \\ \cline{3-5} 
                  &                   &  &  &    \\ \cline{3-5}
                  &                   &TODO  &  &    \\ \hline
                  
                  
                  
\multirow{4}{*}{}8 & \multirow{4}{*}{} TestApproveAccessRequestView& test-not-logged-in & Getestet, ob man ein Zugriffsrequest genehmigen kann, wenn man nicht eingeloggt ist.& HTTP Status Code 302: Der Benutzer wird zur Loginseite weitergeleitet   \\ \cline{3-5}  
                  &                   &  &  &    \\ \cline{3-5}
                  &                   & test-not-owner &Getestet, ob ein Benutzer, der die Ressource nicht besitzt, ein Zugriffsrequest genehmigen kann. & 403  \\ \cline{3-5}
                  &                   & TODO  &  &    \\ \hline
                  
                  
                  
\multirow{4}{*}{} 9& \multirow{4}{*}{} TestDenyAccessRequestView& test-not-logged-in & Getestet, ob man ein Zugriffsrequest verweigern kann, wenn man nicht eingeloggt ist.& HTTP Status Code 302: Der Benutzer wird zur Loginseite weitergeleitet  \\ \cline{3-5}  
                  &                   &  &  &    \\ \cline{3-5}
                  &                   & test-not-owner &Getestet, ob ein Benutzer, der die Ressource nicht besitzt, ein Zugriffsrequest verweigern kann. & 403   \\ \cline{3-5}
                  &                   & TODO  &  &    \\ \hline
                  
                  
                  
\multirow{6}{*}{} 10& \multirow{6}{*}{} TestSendAccessRequestView& test-not-logged-in & Getestet, ob man ein Zugriffsrequest senden kann, wenn man nicht eingeloggt ist.& HTTP Status Code 302: Der Benutzer wird zur Loginseite weitergeleitet \\ \cline{3-5} 
                  &                   &  &  &    \\ \cline{3-5}
				  &                   & test-reader &Getestet, ob ein Benutzer, der die Ressource schon zugreifen darf, ein Zugriffsrequest senden kann. & HTTP Status Code 302: Der Benutzer wird zur Loginseite weitergeleitet  \\ \cline{3-5}
                  &                   & test-staff-user &Getestet, ob der Admin ein Zugriffsrequest senden kann.  &HTTP Status Code 302: Der Benutzer wird zur Loginseite weitergeleitet  \\ \cline{3-5}  
                  &                   &  &  &    \\ \cline{3-5}
                  &                   &TODO  &  &    \\ \hline
                  
                  
                  
\multirow{6}{*}{} 11& \multirow{6}{*}{} TestCancelAccessRequestView& test-not-logged-in & Getestet, ob man ein Zugriffsrequest ablehnen kann, wenn man nicht eingeloggt ist.& HTTP Status Code 302: Der Benutzer wird zur Loginseite weitergeleitet \\ \cline{3-5} 
                  &                   &  &  &    \\ \cline{3-5}
                  &                   & test-reader &Getestet, ob ein Benutzer, der die Ressource schon zugreifen darf, ein Zugriffsrequest ablehnen kann. & HTTP Status Code 302: Der Benutzer wird zur Loginseite weitergeleitet \\ \cline{3-5} 
                  &                   & test-staff-user &Getestet, ob der Admin ein Zugriffsrequest ablehnen kann(Der Admin kann keine Zugriffsrequeste senden ).  &HTTP Status Code 302: Der Benutzer wird zur Loginseite weitergeleitet  \\ \cline{3-5} 
                  &                   &  &  &    \\ \cline{3-5}

                  &                   &TODO  &  &    \\ \hline
                  
                  
                  
\multirow{4}{*}{} 12& \multirow{4}{*}{} TestDeleteResourceView&  &  &   \\ \cline{3-5} 
                  &                   &  &  &    \\ \cline{3-5} 
				  &                   &  &  &    \\ \cline{3-5} 
                  &                   &  &  &    \\ \hline
                  
                  
\multirow{2}{*}{} 13& \multirow{2}{*}{} TestEditNameView&  &  &    \\ \cline{3-5} 
                  &                   &  &  &    \\ \hline
                  
                  
\multirow{3}{*}{} 14& \multirow{3}{*}{} TestResourcesOverview&  &  &    \\ \cline{3-5} 
                  &                   &  &  &    \\ \cline{3-5} 
                  &                   &  &  &    \\ \hline
                  
                  
\multirow{4}{*}{}15 & \multirow{4}{*}{} TestResourcesOverviewSearch&  &  &    \\ \cline{3-5} 
                  &                   &  &  &    \\ \cline{3-5} 
                  &                   &  &  &    \\ \cline{3-5} 
                  &                   &  &  &    \\ \hline
                  
                  
\multirow{8}{*}{}16& \multirow{8}{*}{} TestPermissionEditingView&  &  &    \\ \cline{3-5} 
                  &                   &  &  &    \\ \cline{3-5} 
                  &                   &  &  &    \\ \cline{3-5} 
                  &                   &  &  &    \\ \cline{3-5} 
                  &                   &  &  &    \\ \cline{3-5} 
                  &                   &  &  &    \\ \cline{3-5} 
                  &                   &  &  &    \\ \cline{3-5} 
                  &                   &  &  &    \\ \hline
                  
                  
\multirow{8}{*}{} 17& \multirow{8}{*}{} TestPermissionEditingViewSearch&  &  &   \\ \cline{3-5} 
                  &                   &  &  &    \\ \cline{3-5} 
                  &                   &  &  &    \\ \cline{3-5} 
                  &                   &  &  &    \\ \cline{3-5} 
                  &                   &  &  &    \\ \cline{3-5} 
                  &                   &  &  &    \\ \cline{3-5} 
				  &                   &  &  &    \\ \cline{3-5} 
                  &                   &  &  &    \\ \hline
                  
                  
\multirow{3}{*}{} 18& \multirow{7}{*}{} TestAddNewResourceView&  &  &   \\ \cline{3-5} 
                  &                   &  &  &    \\ \cline{3-5}  
                  &                   &  &  &    \\ \hline
                  
                  
\multirow{4}{*}{} 19& \multirow{7}{*}{} TestOpenResourceView&  &  &   \\ \cline{3-5} 
                  &                   &  &  &    \\ \cline{3-5}  
                  &                   &  &  &    \\ \cline{3-5}  
                  &                   &  &  &    \\ \hline
                  
                  
\multirow{4}{*}{} 20& \multirow{7}{*}{} TestApproveDeletionRequestView&  &  &   \\ \cline{3-5} 
                  &                   &  &  &    \\ \cline{3-5}  
                  &                   &  &  &    \\ \cline{3-5}  
                  &                   &  &  &    \\ \hline
                  
                  
\multirow{4}{*}{} 21& \multirow{7}{*}{} TestDenyDeletionRequestView&  &  &   \\ \cline{3-5} 
                  &                   &  &  &    \\ \cline{3-5}  
                  &                   &  &  &    \\ \cline{3-5}  
                  &                   &  &  &    \\ \hline
                  
                  
\end{longtable}
\newpage


\section{Integration-Test}
Test des Zusammenspiels der Komponenten.

\newpage
\section{System-Test}
System-Test im Sinne von Software Entwicklung ist der abschließende Test, der vom Entwickler durchgeführt wird. Das Ziel ist das Verhalten des Produkts als Ganzes in einer realen Umgebung (immer noch ohne Kunden) zu beobachten.\\
Im Rahmen unseres Projekts sind die System-Tests als manuelle Tests bezeichnet. Sie entsprechen möglichen Szenarien, die durch Interaktion mit dem System entstehen. Für die klare Unterscheidung zwischen den verschiedenen Rollen, die man im Portal haben kann, sind die manuelle Tests in drei Gruppen aufgeteilt - für Benutzer, für Besitzer und für Administrator.

\newpage
\section{Zusammenfassung}
Testüberdeckung. Fazit über die Qualität des Produkts.
\end{document}
\grid
